\section{The Hand-in Archive}
\label{app:handin}

The archive holds this report as \texttt{report.pdf}, a \texttt{README.md} with
the steps below, and two folders, laid out as Table~\ref{tab:handin} shows.
The first, \texttt{task-3-code-generation/}, is an Eclipse project for reading
the transformation and its results, imported with File > Import > Existing
Projects into Workspace. The second, \texttt{implementation/}, contains the EMF
implementation and the examples its tests read.

\begin{table}[!htbp]
\centering
\small
\begin{tabular}{>{\raggedright\arraybackslash}p{0.5\linewidth}>{\raggedright\arraybackslash}p{0.42\linewidth}}
\hline
\textbf{Path} & \textbf{Contents} \\
\hline
\path{model/resolved-deployment.ecore} & The target metamodel, the templates' input. \\
\path{templates/render.mtl} & The Acceleo module. \\
\path{examples/<example>/<project>.resolveddeployment} & The resolved model of \texttt{minimal} and of \texttt{data} the transformation wrote, the templates' source models. \\
\path{generated/} & The tree the templates wrote from both, as the build left it. \\
\path{expected/} & The committed trees it equals: \path{minimal/}, \path{data/} and \path{_estate/}. \\
\path{evidence/} & The output of \texttt{kustomize build} for each directory and of \texttt{kubeconform} for the tree. \\
\hline
\end{tabular}
\caption{The Project of the Archive}
\label{tab:handin}
\end{table}

Every resolved model declares its metamodel in \texttt{xsi:schemaLocation}, so
Eclipse opens it in the Sample Reflective Ecore Model Editor without a
registered package.

\texttt{./mvnw clean verify} in \texttt{implementation/emf} builds the
implementation and runs every test; it needs JDK~21. It parses every example,
runs both transformations and the templates, writes the rendered tree under
\path{bundles/cli/target/parity/rendered/}, and runs the parity test that
compares it with the committed trees. To run the templates in Eclipse Modeling
Tools with the Acceleo~4 SDK, import \texttt{implementation/emf} as existing
Maven projects after that build, and run \texttt{render.launch} in
\texttt{dev.jorisjonkers.deploykit.emf.render}. It renders \texttt{minimal}
from the resolved model the build wrote.
